%! Two Categorical Variables — Unit 2, Lesson 2.1 — Guided Notes & Practice
% THE in-class document, in the MAIN IDEAS / NOTES shape (LESSON_SHAPE.md §1): the vocabulary
% box, then ONE two-column guidednotes table. Each row is one idea — a label on the left; on the
% right a terse definition the student completes, then a grid of short numbered problems with
% reserved answer space. Problems are numbered 1..N through the whole component (14 here).
%   I DO   : the teacher fills the definition lines and works the FIRST problem of each grid
%            (problems 1, 4, 7).
%   WE DO  : the class works the rest of each grid, students holding the pen; the last row,
%            Guided Practice (11–14), is worked entirely together in a fresh context.
%   YOU DO : nothing on this page — ../homework/ IS the individual practice, started in the
%            last minutes of the period.
% COMPACT BUDGET (2026-09-29): 2 pages at 12pt, blank = key.
% THE TRAP is problem 6 (a conditional whose "of those who…" group is the COLUMN, not the row).
% THE CRUX is problem 8 (48 vs 20 raw counts from groups of 160 and 40 — compare percents).
%
% INSTRUCTION DATA (200 students, phone kept in the bedroom overnight × sleep on school nights):
%                 8+ hours   Under 8   Total
%   Phone in room     48        112      160     -> 30% get 8+
%   No phone          20         20       40     -> 50% get 8+
%   Total             68        132      200
% GUIDED PRACTICE DATA (200 drivers, age group × texted while driving in the past month):
%                 Texted   Did not   Total
%   16–24            20        30       50     -> 40%
%   25 and older     36       114      150     -> 24%
%   Total            56       144      200
%
% PAGE LOCKSTEP with notes_key/: every \blank{} here becomes a short \ans{} there, every
% \termblank becomes a \termans, and every \pcell{n}{...}{H}{} gets its answer in the fourth
% argument. Nothing else differs between the two files.
% TABLE MECHANICS: inside a cell, `\\` ENDS THE ROW — break lines with \par. A row cannot break
% across pages: long ideas are split at a seam with a bare `\\` and continued with `& ...`.
% NO SKETCH QUESTIONS: the segmented bar graph is pre-drawn in TikZ and read.
\documentclass[12pt]{article}
\usepackage{apstats-article}
\usepackage{apstats-boxes}

\begin{document}

\pageheader{Unit 2, Lesson 2.1}{Guided Notes \& Practice}
% NAMESTRIP: no \namedateperiod here — the name row lives on the cover only.

\vspace{2pt}

% ── VOCABULARY (I do — students fill each term AS IT IS NAMED, never in advance) ──
\boxguard[16]
\begin{vocabbox}
\small
Fill in each term \textbf{as we name it} in the notes below.
\par
\termblank{Two-way table}
\par
\termblank{Conditional relative frequency}
\par
\termblank{Association}
\par
\end{vocabbox}

\small
\begin{guidednotes}
% ── ROW 1: two-way tables, joint and marginal ────────────────────────────────
\mainidea[Reading a]{Two-way table} &
\begin{minipage}[c]{0.40\linewidth}\raggedright
200 students: does your phone stay in your bedroom overnight? Do you sleep 8+ hours on a
school night?
\end{minipage}\hfill
\begin{minipage}[c]{0.58\linewidth}
\centering\footnotesize
\setlength{\tabcolsep}{3pt}\renewcommand{\arraystretch}{1.15}
\begin{tabular}{@{} l | c c | c @{}}
 & \textbf{8+ hrs} & \textbf{Under 8} & \textbf{Total} \\ \hline
\textbf{Phone in room} & 48 & 112 & 160 \\
\textbf{No phone} & 20 & 20 & 40 \\ \hline
\textbf{Total} & 68 & 132 & 200 \\
\end{tabular}
\end{minipage}

\vspace{4pt}
A \textbf{two-way table} summarizes two \blank{2.3cm} variables on the same individuals. A
\textbf{joint} relative frequency is one cell $\div$ the \blank{1.9cm} total. A
\textbf{marginal} relative frequency is a row or column total $\div$ the \blank{1.9cm} total.
\notesprompt{Answer as a proportion of all 200 students.}
\begin{probgrid}{|Y|Y|Y|}
\pcell{1}{Phone in room \emph{and} 8+ hours}{0.9cm}{} &
\pcell{2}{Students who sleep 8+ hours}{0.9cm}{} &
\pcell{3}{``0.80 keep a phone in the room'' --- joint or marginal?}{0.9cm}{} \\ \hline
\end{probgrid}
\\ \hline
% ── ROW 2: conditional relative frequency — pick the denominator ─────────────
\mainidea[Of those who\ldots]{Conditional} &
A \textbf{conditional} relative frequency looks inside \blank{1.1cm} row or column only: divide by
that group's \blank{1.3cm}. In ``of those who \ldots'', the words after \emph{of} name the
\blank{2.2cm}.
\notesprompt{Find each conditional relative frequency. Show the fraction.}
\begin{probgrid}{|Y|Y|Y|}
\pcell{4}{Of phone-in-room students: proportion with 8+ hours?}{1.0cm}{} &
\pcell{5}{Of no-phone students: proportion with 8+ hours?}{1.0cm}{} &
\pcell{6}{Of 8+ hour sleepers: proportion with a phone in the room?}{1.0cm}{} \\ \hline
\end{probgrid}
\\ \hline
% ── ROW 3: association from a segmented bar graph — THE CRUX ─────────────────
\mainidea[Comparing groups:]{Association} &
\begin{minipage}[c]{0.40\linewidth}
\centering
\begin{tikzpicture}[x=1cm, y=0.03cm]
  \foreach \y in {0,50,100} {
    \draw (-0.08,\y) -- (0.08,\y) node[left, font=\tiny, xshift=-2pt] {\y\%};}
  \fill[navy]    (0.3,0)  rectangle (1.3,30);
  \fill[goldacc] (0.3,30) rectangle (1.3,100);
  \fill[navy]    (1.8,0)  rectangle (2.8,50);
  \fill[goldacc] (1.8,50) rectangle (2.8,100);
  \draw (0.3,0) rectangle (1.3,100);
  \draw (1.8,0) rectangle (2.8,100);
  \node[white, font=\scriptsize\bfseries] at (0.8,15) {30\%};
  \node[font=\scriptsize\bfseries] at (0.8,65) {70\%};
  \node[white, font=\scriptsize\bfseries] at (2.3,25) {50\%};
  \node[font=\scriptsize\bfseries] at (2.3,75) {50\%};
  \draw[thick] (0,0) -- (0,104);
  \draw[thick] (0,0) -- (3.0,0);
  \node[below, font=\tiny, align=center] at (0.8,-2) {Phone in\\room (160)};
  \node[below, font=\tiny, align=center] at (2.3,-2) {No phone\\(40)};
  \fill[navy] (0.3,112) rectangle (0.5,120); \node[right, font=\tiny] at (0.5,116) {8+ hours};
  \fill[goldacc] (1.8,112) rectangle (2.0,120); \node[right, font=\tiny] at (2.0,116) {Under 8};
\end{tikzpicture}
\end{minipage}\hfill
\begin{minipage}[c]{0.57\linewidth}\raggedright
A \textbf{segmented bar graph} draws each group's conditional distribution as one bar that
totals \blank{1.3cm}. (A \textbf{mosaic plot} also makes each bar's width show the group's
\blank{1.1cm}.)\par\vspace{3pt}
Two variables are \textbf{associated} when the conditional distributions \blank{1.6cm} from
group to group. Groups of different sizes: compare \blank{1.7cm}, not counts.
\end{minipage}
\\
& \notesprompt{Use the graph and the table above.}
\begin{probgrid}{|Y|Y|}
\pcell{7}{What percent of phone-in-room students sleep under 8 hours?}{0.9cm}{} &
\pcell{8}{A student says: ``48 phone-in-room students sleep 8+ hours, but only 20 no-phone
students do --- phones go with \emph{more} sleep.'' What is wrong?}{1.6cm}{} \\ \hline
\pcell{9}{Are phone location and sleep associated? Justify with percents.}{1.6cm}{} &
\pcell{10}{Does this survey show that a phone in the room \emph{causes} less sleep? Why?}{1.6cm}{} \\ \hline
\end{probgrid}
\\ \hline
% ── LAST ROW: GUIDED PRACTICE — we do, together, a fresh context ─────────────
\mainidea[Guided practice]{Texting drivers} &
\begin{minipage}[c]{0.40\linewidth}\raggedright
A county survey asked 200 drivers their age and whether they had texted while driving in the
past month.
\end{minipage}\hfill
\begin{minipage}[c]{0.58\linewidth}
\centering\footnotesize
\setlength{\tabcolsep}{3pt}\renewcommand{\arraystretch}{1.15}
\begin{tabular}{@{} l | c c | c @{}}
 & \textbf{Texted} & \textbf{Did not} & \textbf{Total} \\ \hline
\textbf{Age 16--24} & 20 & 30 & 50 \\
\textbf{Age 25+} & 36 & 114 & 150 \\ \hline
\textbf{Total} & 56 & 144 & 200 \\
\end{tabular}
\end{minipage}
\notesprompt{We work these together.}
\begin{probgrid}{|Y|Y|}
\pcell{11}{What proportion of all 200 drivers texted? Joint or marginal?}{1.2cm}{} &
\pcell{12}{Find the proportion who texted in each age group.}{1.2cm}{} \\ \hline
\pcell{13}{A headline reads: ``Older drivers text more --- 36 admitted it, versus 20 young
drivers.'' Respond.}{1.7cm}{} &
\pcell{14}{Are age and texting associated in this survey? Answer in context.}{1.7cm}{} \\ \hline
\end{probgrid}
\\ \hline
\end{guidednotes}

% ── THE NOTES END AT GUIDED PRACTICE ─────────────────────────────────────────
% There is deliberately no "you do" here: ../homework/ is the individual practice,
% started in class in the last 8 minutes while the teacher circulates.

\end{document}
